\documentclass{article}
\usepackage[a4paper, margin=.7in]{geometry}
\usepackage{amsmath}
\usepackage{amssymb}
\usepackage{hyperref}
\usepackage{nameref}
\usepackage{algorithm}
\usepackage{algpseudocode}

% No section numbering
\setcounter{secnumdepth}{0}
% No paragraph indentation
\setlength{\parindent}{0pt}

% Define custom commands
%% Bold-italic math symbols
\newcommand{\mbi}[1]{\boldsymbol{#1}}
%% Manually indented state in algorithmic environment
\newcommand{\IndentState}{\State \hspace{1em}}
\newcommand{\DIndentState}{\State \hspace{2em}}
%% Argmin operator
\DeclareMathOperator*{\argmin}{arg\,min}
%% Argmax operator
\DeclareMathOperator*{\argmax}{arg\,max}

\title{CSCI-567 Assignment 3 Solutions}
\author{
  Naman Chhibbar \\
  USC ID: 7709-8548-02
}
\date{}

% Document body
\begin{document}
\maketitle

\section{Problem 1}
  \begin{itemize}
    \item[\textbf{1.1}]
    \begin{itemize}
      \item[(1)] The possible outcomes are $A$ and $C$, and $B$ and $C$.
      This is because if one of the centers is at $A$ or $B$, then the other one must be at $C$ due to greedy initialization.

      \item[(2)] Due to symmetry, we can safely assume that the centers of clusters 1 and 2 are initialized at $A$ and $C$, respectively.
      K-means will converge in a single iteration with points at $A$ and $B$ in cluster 1, and the only point at $C$ in cluster 2.
      Since we have equal number of points on $A$ and $B$, the center of cluster 1 will be at point $D$ (midpoint of $A$ and $B$).
      Since cluster 2 has only one point, its center will remain at $C$.

      \item[(3)] We can see that all $2m$ points in cluster 1 are at a distance of 1 from the center.
      Whereas, the single point in cluster 2 is at 0 distance.
      Hence, the K-means objective is $2m$.
    \end{itemize}

    \item[\textbf{1.2}]
    \begin{itemize}
      \item[(1)] Let $p$ and $p'$ be the probabilities of choosing $A$ and $B$, and $A$ and $C$ as the initial centers, respectively.
      By symmetry, the probability of choosing $B$ and $C$ is also $p'$.
      Hence, we have $p + 2p' = 1$.
      Now:
      \[ p = \mathbb{P}(\text{choosing }A) \cdot \mathbb{P}(\text{choosing }B | A \text{ is chosen}) \]
      Since the first center is randomly chosen, we have:
      \[ \mathbb{P}(\text{choosing }A) = \frac{m}{2m + 1} \]
      because there are $m$ points at $A$ out of a total of $(2m + 1)$ points.
      Now, if $A$ is the first chosen center, then the probability of choosing a point is proportional to its distance from $A$ squared.
      There are $m$ points at $B$ each at a distance of 2 from $A$, and the only point at $C$ is at a distance of $\sqrt{10}$ from $A$.
      Hence:
      \[ \mathbb{P}(\text{choosing }B | A \text{ is chosen}) = \frac{m \cdot 2^2}{m \cdot 2^2 + 1 \cdot (\sqrt{10})^2} = \frac{4m}{4m + 10} = \frac{2m}{2m + 5} \]
      Now, we can calculate $p$ as:
      \[ p = \frac{m}{2m + 1} \cdot \frac{2m}{2m + 5} = \frac{2m^2}{(2m + 1)(2m + 5)} \]
      We also get $p' = (1-p)/2$.

      \item[(2)] In such a case, every point is at their respective centers, except for $x_{2m+1}$.
      Whether $x_{2m+1}$ is assigned to cluster 1 or 2, the K-means objective will be 10 (since $C$ is at a distance of $\sqrt{10}$ from both $A$ and $B$).

      \item[(3)] We have seen that if $A$ and $B$ (we assume that $m$ is large enough for this) are the initial centers then the K-means objective is 10 (in part (2)), and if $A$ and $C$ (or $B$ and $C$) are the initial centers then the K-means objective is $2m$ (in part (3) of question \textbf{1.1}).
      Hence, the expected K-means objective is:
      \begin{align*}
        \mathbb{E}[\text{K-means objective}] &= p \cdot 10 + 2p' \cdot 2m \\
        &= 10p + 4m \cdot \frac{1-p}{2} \\
        &= (10 - 2m)p + 2m \\
      \end{align*}
      where we have used $p + 2p' = 1$ and $p = \frac{2m^2}{(2m + 1)(2m + 5)}$.
      Also, note that:
      \begin{align*}
        \lim_{m \to \infty} p &= \lim_{m \to \infty} \frac{2m^2}{(2m + 1)(2m + 5)} \\
        &= \lim_{m \to \infty} \frac{2}{(2 + 1/m)(2 + 5/m)} \\
        &= \frac{1}{2}
      \end{align*}
      Hence:
      \begin{align*}
        \lim_{m \to \infty} \mathbb{E}[\text{K-means objective}] &= \lim_{m \to \infty} \left( (10 - 2m)p + 2m \right) \\
        &= \lim_{m \to \infty} (10 - 2m) \cdot \frac{1}{2} + 2m \\
        &= \lim_{m \to \infty} 5 - m + 2m \\
        &= \lim_{m \to \infty} 5 + m \\
        &= \infty
      \end{align*}
    \end{itemize}
    
    \item[\textbf{1.3}] In question \textbf{1.2}, we saw that as $m \to \infty$, $p \to 1/2$ and hence $p' \to 1/4$.
    This means that choosing $A$ and $B$ are likely to be chosen as the initial centers when $m$ is large, showing the robustness to outliers of this initialization.
    Whereas, greedy initialization will always choose either $A$ and $C$, or $B$ and $C$, leading to a local minimum.
    Even though the expectation of K-means objective goes to infinity as $m \to \infty$ when using K-means++ initialization, the probability of getting a good initialization is still higher.
  \end{itemize}
  
\section{Problem 2}
  \begin{itemize}
    \item[\textbf{2.1}] By Bayes' theorem, we have:
    \begin{align*}
      q_n(z) &= p(z_n = z | \mbi{x}_n; \mbi{v}) \\
      &= \frac{p(z_n = z; \mbi{v})}{p(\mbi{x}_n; \mbi{v})} \cdot p(\mbi{x}_n | z_n = z; \mbi{v}) \\
      &\propto p(z_n = z) \cdot p(\mbi{x}_n | z_n = z; \mbi{v}) \\
      &\propto \exp\left(-\frac{1}{2} z^2\right) \cdot \exp\left(-\frac{1}{2\sigma^2} ||\mbi{x}_n - z \mbi{v}||^2\right) \\
      &\propto \exp\left(-\frac{z^2}{2} - \frac{||\mbi{x}_n||^2 + z^2||\mbi{v}||^2 - 2z \mbi{x}_n^T \mbi{v}}{2\sigma^2}\right) \\
      &\propto \exp\left(-\frac{(\sigma^2 + ||\mbi{v}||^2)z^2 - 2z \mbi{x}_n^T \mbi{v}}{2\sigma^2}\right) \\
      &\propto \exp\left(-\frac{1}{2 \cdot \frac{\sigma^2}{\sigma^2 + ||\mbi{v}||^2}} \cdot \left(z^2 - 2z \frac{\mbi{x}_n^T \mbi{v}}{\sigma^2 + ||\mbi{v}||^2}\right)\right) \\
      &\propto \exp\left(-\frac{1}{2 \cdot \frac{\sigma^2}{\sigma^2 + ||\mbi{v}||^2}} \cdot \left(z - \frac{\mbi{x}_n^T \mbi{v}}{\sigma^2 + ||\mbi{v}||^2}\right)^2\right) \\
      &\propto \exp\left(-\frac{1}{2b^2} \cdot (z - a_n)^2\right)
    \end{align*}
    where $a_n = \frac{\mbi{x}_n^T \mbi{v}}{\sigma^2 + ||\mbi{v}||^2}$ and $b^2 = \frac{\sigma^2}{\sigma^2 + ||\mbi{v}||^2}$.
    Hence, the posterior distribution is a one-dimensional Gaussian with mean $a_n$ and variance $b^2$.
  \end{itemize}
  
\section{Problem 3}
  \begin{itemize}
    \item[\textbf{3.1}]
    \begin{itemize}
      \item[1)] Let $n \in \{1, 2, \ldots, N\}$.
      Maximizing equation (1) with respect to $c_n$:
      \begin{align*}
        \frac{\partial}{\partial c_n} \sum_{i=1}^N ||\mbi{x}_i - c_i \mbi{v}||_2^2 &= 0 \\
        \Rightarrow 2(\mbi{x}_n - c_n \mbi{v})^T (-\mbi{v}) &= 0 \\
        \Rightarrow c_n ||\mbi{v}||_2^2 &= \mbi{x}_n^T \mbi{v} \\
        \Rightarrow c_n &= \mbi{x}_n^T \mbi{v} \quad \text{($v$ is a unit vector)}
      \end{align*}
      Therefore, $c_n = \mbi{x}_n^T \mbi{v} ~ \forall n \in \{1, 2, \ldots, N\}$.

      \item[2)]
      \begin{align*}
        \argmin_{\mbi{v}:||\mbi{v}||_2=1} \sum_{n=1}^N ||\mbi{x}_n - c_n \mbi{v}||_2^2 &= \argmin_{\mbi{v}:||\mbi{v}||_2=1} \sum_{n=1}^N ||\mbi{x}_n - (\mbi{x}_n^T \mbi{v}) \mbi{v}||_2^2 \\
        &= \argmin_{\mbi{v}:||\mbi{v}||_2=1} \sum_{n=1}^N ||\mbi{x}_n||_2^2 - 2(\mbi{x}_n^T \mbi{v})^2 + (\mbi{x}_n^T \mbi{v})^2 ||\mbi{v}||_2^2 \\
        &= \argmin_{\mbi{v}:||\mbi{v}||_2=1} \sum_{n=1}^N ||\mbi{x}_n||_2^2 - (\mbi{x}_n^T \mbi{v})^2 \\
        &= \argmin_{\mbi{v}:||\mbi{v}||_2=1} -\sum_{n=1}^N (\mbi{x}_n^T \mbi{v})^2 \quad \text{($||\mbi{x}_n||_2^2$ does not depend on $\mbi{v}$)} \\
        &= \argmax_{\mbi{v}:||\mbi{v}||_2=1} \sum_{n=1}^N (\mbi{x}_n^T \mbi{v})^2 \\
        &= \argmax_{\mbi{v}:||\mbi{v}||_2=1} \sum_{n=1}^N \mbi{v}^T \mbi{x}_n \mbi{x}_n^T \mbi{v} \\
        &= \argmax_{\mbi{v}:||\mbi{v}||_2=1} \mbi{v}^T \left( \sum_{n=1}^N \mbi{x}_n \mbi{x}_n^T \right) \mbi{v} \\
      \end{align*}
      Hence, by solving equation (1), we get the first principal component of the dataset.
    \end{itemize}

    \item[\textbf{3.2}]
    \begin{itemize}
      \item[1)] The reconstructed dataset is $\mathbf{V}\mbi{c}_1, \ldots, \mathbf{V}\mbi{c}_N$.
      Note that each $\mathbf{V}\mbi{c}_n \in \mathbb{R}^D$.
      Also, the analogous optimization problem is:
      \[ \argmin_{\mbi{c}_1, \ldots, \mbi{c}_N, \mathbf{V}: \mathbf{V}^T \mathbf{V} = \mathbf{I}} \sum_{n=1}^N ||\mbi{x}_n - \mathbf{V}\mbi{c}_n||_2^2 \]

      \item[2)] Let $n \in \{1, 2, \ldots, N\}$.
      Maximizing the above equation with respect to $\mbi{c}_n$:
      \begin{align*}
        \frac{\partial}{\partial \mbi{c}_n} \sum_{i=1}^N ||\mbi{x}_i - \mathbf{V}\mbi{c}_i||_2^2 &= 0 \\
        \Rightarrow -2 \mathbf{V}^T (\mbi{x}_n - \mathbf{V}\mbi{c}_n) &= 0 \\
        \Rightarrow \mathbf{V}^T \mathbf{V} \mbi{c}_n &= \mathbf{V}^T \mbi{x}_n \\
        \Rightarrow \mbi{c}_n &= \mathbf{V}^T \mbi{x}_n \quad \text{(since columns of $\mathbf{V}$ are orthonormal)}
      \end{align*}
      Therefore, $\mbi{c}_n = \mathbf{V}^T \mbi{x}_n ~ \forall n \in \{1, 2, \ldots, N\}$.

      \item[3)] Plugging the value of $\mbi{c}_n$ back into the optimization problem:
      \begin{align*}
        \argmin_{\mathbf{V}: \mathbf{V}^T \mathbf{V} = \mathbf{I}} \sum_{n=1}^N ||\mbi{x}_n - \mathbf{V} \mathbf{V}^T \mbi{x}_n||_2^2 &= \argmin_{\mathbf{V}: \mathbf{V}^T \mathbf{V} = \mathbf{I}} \sum_{n=1}^N ||\mbi{x}_n||_2^2 - 2 \mbi{x}_n^T \mathbf{V} \mathbf{V}^T \mbi{x}_n + ||\mathbf{V} \mathbf{V}^T \mbi{x}_n||_2^2 \\
        &= \argmin_{\mathbf{V}: \mathbf{V}^T \mathbf{V} = \mathbf{I}} \sum_{n=1}^N ||\mbi{x}_n||_2^2 - ||\mathbf{V}^T \mbi{x}_n||_2^2 \quad \text{(since columns of $\mathbf{V}$ are orthonormal)} \\
        &= \argmin_{\mathbf{V}: \mathbf{V}^T \mathbf{V} = \mathbf{I}} - \sum_{n=1}^N ||\mathbf{V}^T \mbi{x}_n||_2^2 \quad \text{($||\mbi{x}_n||_2^2$ does not depend on $\mathbf{V}$)} \\
        &= \argmax_{\mathbf{V}: \mathbf{V}^T \mathbf{V} = \mathbf{I}} \sum_{n=1}^N ||\mathbf{V}^T \mbi{x}_n||_2^2 \\
        &= \argmax_{\mathbf{V}: \mathbf{V}^T \mathbf{V} = \mathbf{I}} \sum_{n=1}^N \mbi{x}_n^T \mathbf{V} \mathbf{V}^T \mbi{x}_n \\
        &= \argmax_{\mathbf{V}: \mathbf{V}^T \mathbf{V} = \mathbf{I}} \sum_{j=1}^p \left( \mbi{v}_j^T \left( \sum_{n=1}^N \mbi{x}_n \mbi{x}_n^T \right) \mbi{v}_j \right) \\
      \end{align*}
    \end{itemize}
  \end{itemize}

\end{document}
